\section{Anwendungsfall mit Messdaten der GeoSphere Austria}
\label{sec:anwendungsfall}

% VERANTWORTLICH: ...
% Aufgabenpunkt 5, Zeitraum 24.06. bis 01.07.2019.

\subsection{Datengrundlage}

% Station, Zeitaufloesung, Zeitzone, Umgang mit Luecken.
% Falls Horizontalstrahlung verwendet wird und nicht die Strahlung in
% Modulebene: das ist eine Annahme und gehoert in Tabelle 2.1.

\subsection{Ergebnisse}

\begin{figure}[htbp]
  \centering
  % \includegraphics[width=\textwidth]{anwendungsfall_bilanz}
  \caption{Zeitlicher Verlauf der einzelnen Terme der Waermebilanz nach
           \cref{eq:bilanz}. Der konvektive Verlust dominiert waehrend der
           Tagesstunden, der Speicherterm bleibt klein.}
  \label{fig:bilanz}
\end{figure}

\subsection{Dominanter Verlustterm}

% Die Angabe fragt explizit: welcher Verlustterm spielt die groesste Rolle
% und was folgt daraus fuer einen effizienteren Betrieb?
% Antwort quantitativ belegen, ueber die integrierten Energiemengen, nicht
% nur ueber den Blick auf die Kurve.
% Daraus ableiten: Massnahmen zur besseren Kuehlung, Hinterlueftung,
% Montageart, und was das an Mehrertrag bringen wuerde.
